\introduction

\section*{Context and Motivation}
\addcontentsline{toc}{section}{Context and Motivation}

Planning vehicle routes is an important task in logistics. A business must
decide how to use its fleet to serve customers while respecting operational
requirements and controlling costs. Vehicle routing problems provide a
framework for studying these decisions~\cite{toth2014vehicle}. When a vehicle
can perform several trips, planning also involves deciding how those trips
should follow one another~\cite{cattaruzza2016multipletrips}.

For vehicles transporting different products, the cost of a plan can depend
on more than the distance travelled. Changing from one product to another
may require cleaning or preparation before the next load. For example, in
an illustrative bulk-delivery operation, a vehicle may need cleaning before
carrying a product that is incompatible with residues from its previous
load. The required preparation can depend on both products and on their
order. A plan with shorter journeys can therefore incur higher overall
costs if it requires more expensive product changes.

This observation motivates the study of routing decisions together with
\emph{changeover costs}: the costs associated with preparing a vehicle to
carry a product after the one previously transported.

\section*{Problem Statement}
\addcontentsline{toc}{section}{Problem Statement}

Cleaning and product-related requirements have already been considered in
routing research, including the collection of different grades of olive
oil~\cite{lahyani2015oilmulticompartment}. These studies provide relevant
precedents. The present work examines a multi-product setting in which
vehicles perform repeated loading and delivery operations, and the costs of
successive product changes are included in the planning objective. This
problem is referred to as the \emph{Multi-Product Vehicle Routing Problem
with Changeover Costs} (MPVRP-CC). The challenge is to coordinate deliveries
and product sequences while respecting the operational constraints and
controlling both travel and changeover costs.

This thesis investigates constraint programming (CP) as an approach to this
problem. CP allows operational requirements to be expressed as rules that
solutions must satisfy. In particular, sequence variables provide a way to
represent and reason about ordered operations in routing
problems~\cite{delecluse2022sequence}. Their use raises both a modeling
question and an experimental one: how to represent MPVRP-CC, and how well the
resulting solution approach performs.

The central research question is:
\begin{quote}
\emph{To what extent can constraint programming with sequence variables
provide effective solutions to MPVRP-CC, and how does accounting for
changeover costs influence route planning and total modeled operational cost?}
\end{quote}
Three specific questions guide the study:
\begin{enumerate}
  \item How can the operational constraints and changeover costs of MPVRP-CC
        be represented in a solution approach using CP and sequence variables?
  \item How does the proposed approach compare with a mixed-integer linear
        programming (MILP) baseline in terms of feasible solutions found,
        solution cost, and computation time?
  \item How do the resulting vehicle routes, product sequences, and total
        modeled costs differ when changeover costs are included in planning
        rather than ignored?
\end{enumerate}

\section*{Objectives}
\addcontentsline{toc}{section}{Objectives}

The main objective is to develop and evaluate an approach to MPVRP-CC using
constraint programming and sequence variables, and to examine the influence
of changeover costs on planning decisions. The specific objectives are to:
\begin{enumerate}
  \item define MPVRP-CC, its operational constraints, and its cost objective,
        and formulate a MILP reference model;
  \item design and implement a solution approach using CP and sequence
        variables;
  \item assess the approach against the MILP baseline on synthetic instances;
        and
  \item analyze the effect of including changeover costs on the resulting
        plans and their total modeled costs.
\end{enumerate}

The study uses 100 synthetic instances with varied configurations. Its
conclusions concern these instances and the operational assumptions of the
model. The methodological choices developed during the research, including
the construction of delivery groups before CP assignment and sequencing,
are explained in the Methodology chapter.

\section*{Thesis Organization}
\addcontentsline{toc}{section}{Thesis Organization}

Chapter~\ref{chap:literature-review} reviews relevant routing variants,
cleaning and changeover costs, and solution methods, and positions this
study within the literature. Chapter~\ref{chap:theoretical-background}
introduces the routing, optimization, and CP concepts needed to understand
the approach. Chapter~\ref{chap:methodology} defines MPVRP-CC, presents the
MILP model, describes the proposed approach, and specifies the experimental
setup. Chapter~\ref{chap:results-discussion} reports the method comparison,
examines the effect of changeover costs, and discusses the findings and
limitations. Chapter~\ref{chap:conclusion-future-work} summarizes the
contributions and outlines future work. The references and appendices
provide supporting sources, data formats, algorithms, visualizations, and
supplementary details.
